\section{Limitations and Differences from the Original Paper}

This work adapts the paper's prediction mechanisms to gem5 MinorCPU.
It does not reproduce the paper's full processor or its evaluation
setup~\cite{ravenel2025architecting}. The two differ in the
processor model, in when predictions are checked, and in how
performance is measured.

\subsection{Processor Model and Evaluation Scope}

The paper evaluates a modified CVA6-based processor on SPEC CPU2017
using SimPoints. This work uses gem5 MinorCPU and small targeted
RISC-V microbenchmarks. Both processors are four-wide. The paper's
processor can issue two loads per cycle and has a 32-entry
scoreboard. The MinorCPU setup used here issues at most one memory
instruction per cycle, and its scoreboard tracks every architectural
register~\cite{ravenel2025architecting}. The caches and workloads
also differ. So the cycle reductions reported here describe this
simulator setup and these microbenchmarks only. They are not
application-level speedups and should not be compared directly with
the paper's results.

\subsection{Differences from the Original Design}
\label{subsec:original-design-differences}

\paragraph{LVP and recovery:} The LVP table and its confidence policy
follow the paper~\cite{ravenel2025architecting}. The difference is in
how predictions are used. Here, an LVP prediction can only serve as
the base address of a dependent integer load, and only while its
producer is the oldest instruction and still waiting for its data
(Section~\ref{sec:lvp-integration}). Both the paper and this work
recover only when a wrong prediction was actually used, but the
recovery method differs. The paper replays younger instructions
through its scoreboard, while the implementation squashes them and fetches
them again.

\paragraph{Address and tag compression:} The paper's Last Address
Compressed Table (LACT) is a 32-entry LRU structure that replaces the
upper 44 bits of an address with a 5-bit identifier and keeps only the
lower 20 bits in the predictor. This cuts predictor storage by 85\%.
The value-cache tags reuse the same identifier, which cuts tag
storage by 75\%~\cite{ravenel2025architecting}. This work keeps full
64-bit bases and full page-number tags instead. The value cache has
the same 4\,KiB capacity and 16-byte entries as the paper's, but it
tracks validity per byte, does not fill whole 16-byte lines.

\paragraph{When V-ASTRA is checked:} The paper validates the predicted
base at issue, so dependents can issue in the same cycle as the load.
Here, the base is checked in Execute when the load forms its address,
and the bypass is accepted later in the LSQ, after address translation
and some of the safety checks (Section~\ref{sec:impl-vcache}).

\paragraph{Store safety:} The paper tracks possible conflicts with
older stores using a counting Bloom filter and an invalidation
FIFO~\cite{ravenel2025architecting}. This work scans the older
in-flight stores exactly, both when forming a candidate and again
before the bypass. An exact scan avoids the Bloom filter's false
positives, but its cost grows with the number of in-flight stores.
The FIFO is not needed here. Memory operations enter the LSQ in
order, so older store addresses are already known at the decision
point.

\paragraph{L0 minicache:} The paper's L0 minicache is a 32-entry
direct-mapped structure with speculative and committed data. It lets
loads blocked by in-flight stores still be served
quickly~\cite{ravenel2025architecting}. It is not implemented here.
Such loads fall back to the normal memory path, where MinorCPU's
store buffer forwards fully covered loads and makes partly covered
ones wait.

\paragraph{Dependency handling:} In the paper, dependents of a
validated load can issue and execute in the same cycle. In MinorCPU,
non-memory instructions still execute at commit. Design~A only lets
dependents issue earlier under stall-on-use timing, and it has no
effect under the default timing. Design~B forwards a validated value,
but only to the base address of a dependent integer load, not to ALU,
branch, store, or floating-point instructions.

\subsection{Future Work}

\paragraph{Extending this work:} The random-base result suggests a
smarter rule for when an LVP-supplied base may be used for V-ASTRA,
for example based on LVP confidence, recent validation results, or
agreement with ASTRA. Such a rule should be tested on random base,
stable flip, V-cache flip, mixed, and the forced-misprediction
workload, to check both gains and regressions. Design~B should be
measured on its own, by comparing V-ASTRA with and without
forwarding, and Design~A should be measured under stall-on-use
timing. Design~B could also forward to other load operands and to
store addresses. Forwarding to ALU and branch instructions would need
larger changes to MinorCPU.

\paragraph{Closer to the paper:} The biggest missing piece is the
L0 minicache, which would give a fast path to loads that V-ASTRA
currently rejects because of in-flight stores. It would need LSQ
integration, speculative store handling, line locking, and recovery,
and it should be evaluated on store-heavy workloads together with
META. Adding LACT-style compression
(Section~\ref{subsec:original-design-differences}) would mainly reduce
storage. It would need identifier allocation and recycling,
compressed tables and tags, cache invalidation when identifiers are
reused, and a check for displacements that cross the compressed
region. Replacing the exact store scan with a Bloom filter, or
squash-and-refetch with scoreboard replay, would bring the design
closer to the paper but is lower priority. The paper also discusses
VTAGE, process-ID-tagged value-cache entries, multiprocessor support,
and security~\cite{ravenel2025architecting}.

\paragraph{Broader evaluation:} The current results use deterministic
microbenchmarks and one processor and memory configuration. Future
evaluations should use application-level and embedded benchmarks,
other cache and memory-latency settings, and other processor widths
where possible. Supporting compressed RISC-V loads would allow
binaries built for \texttt{rv64gc}.
